\documentclass[12pt,a4paper]{article}
\usepackage[margin=25mm, headheight=15pt, headsep=10pt]{geometry}
\usepackage{fontspec}
\usepackage[table]{xcolor} % Note: table option is required for rowcolors
\usepackage{titlesec}
\usepackage{tocloft}
\usepackage{fancyhdr}
\usepackage{booktabs}
\usepackage{tcolorbox}
\tcbuselibrary{skins, breakable}
\usepackage[colorlinks=true, linkcolor=emerald, urlcolor=emerald, bookmarks=true]{hyperref}

\definecolor{emerald}{RGB}{0,102,51}
\definecolor{darkred}{RGB}{139,0,0}
\definecolor{lightgray}{RGB}{245,245,245}

\usepackage[nil]{babel}
\babelprovide[import, main]{bengali}
\babelprovide[import]{arabic}
\babelfont{rm}[Renderer=HarfBuzz,Scale=1.0]{Noto Serif Bengali}
\babelfont{sf}[Renderer=HarfBuzz]{Noto Sans Bengali}
\babelfont{tt}[Renderer=HarfBuzz]{Noto Sans Bengali}
\babelfont[arabic]{rm}[Renderer=HarfBuzz,Scale=1.1]{Noto Naskh Arabic}

% Section styling
\titleformat{\section}{\Large\bfseries\color{emerald}}{\thesection.}{0.5em}{}
\titleformat{\subsection}{\large\bfseries\color{emerald!85!black}}{\thesubsection.}{0.5em}{}
\titleformat{\subsubsection}{\normalsize\bfseries\color{emerald!70!black}}{\thesubsubsection.}{0.5em}{}
\titlespacing*{\section}{0pt}{18pt}{8pt}

% Fancy Header/Footer setup
\pagestyle{fancy}
\fancyhf{} % clear all header and footer fields
\fancyhead[L]{\color{emerald}\small\textbf{\nouppercase{\leftmark}}}
\fancyfoot[L]{\scriptsize\color{black!50}{\lat{AI}}-সংকলিত, আলেম-যাচাইকৃত নয়}
\fancyfoot[C]{\color{emerald!70!black}\thepage}
\renewcommand{\headrulewidth}{0.4pt}
\renewcommand{\headrule}{\hbox to\headwidth{\color{emerald}\leaders\hrule height \headrulewidth\hfill}}
\renewcommand{\footrulewidth}{0pt}

% Custom boxes
% 1. Ayat/Hadith Box
\newtcolorbox{ayatbox}[1][]{
    enhanced, breakable, 
    colback=emerald!5, 
    colframe=emerald, 
    boxrule=0.5pt, 
    arc=3pt, 
    left=10pt, right=10pt, top=10pt, bottom=10pt,
    #1
}

% 2. Quote/Translation Box
\newtcolorbox{quotebox}[1][]{
    enhanced, breakable,
    frame hidden,
    colback=lightgray,
    borderline west={3pt}{0pt}{emerald},
    left=15pt, right=10pt, top=10pt, bottom=10pt,
    sharp corners,
    #1
}

% 3. AI-generated notice box
\newtcolorbox{warnbox}[1][]{
    enhanced, breakable,
    colback=red!4,
    colframe=darkred,
    boxrule=0.8pt,
    arc=3pt,
    left=10pt, right=10pt, top=10pt, bottom=10pt,
    #1
}

% Commands
\newcommand{\arab}[1]{\foreignlanguage{arabic}{#1}}
\newcommand{\kib}[1]{\textcolor{emerald}{\textbf{#1}}}

% Latin (English) fallback: Noto Serif Bengali has NO Latin glyphs
\newfontfamily\latfont[Renderer=HarfBuzz]{Noto Serif}
\newcommand{\lat}[1]{{\latfont #1}}

\setlength{\parskip}{5pt}
\setlength{\parindent}{0pt}

% Fix for missing bullet in Noto Serif Bengali
\renewcommand{\labelitemi}{$\bullet$}



\begin{document}
\begin{center}{\Large\bfseries\color{emerald} 12-উমার-ঘাসের-মুঠা-ক্ষমতা-ও-বিনয়}\end{center}
\vspace{1em}
\section*{ঘটনা ১২: উমার (রাঃ) — মাথার উপর ঘাসের মুঠা}

\begin{warnbox}
 \textbf{গুরুত্বপূর্ণ নোটিশ (\lat{Important} \lat{Notice}):} এই কন্টেন্টটি \lat{AI} (কৃত্রিম বুদ্ধিমত্তা) দিয়ে প্রস্তুত ও সংকলিত। \lat{AI} কঠোর সূত্র-মান নীতি মেনে শুধু নির্ভরযোগ্য উৎস থেকে তথ্য সংগ্রহ করেছে — কেবল সহীহ/হাসান হাদিস ও সহীহ সনদের আসার রাখা হয়েছে, দুর্বল সনদ বাদ দেওয়া হয়েছে। তবে এটি কোনো আলেম বা মুহাদ্দিস কর্তৃক যাচাইকৃত নয়।
\end{warnbox}


\medskip\hrule\medskip

\subsection*{১. সূত্র কার্ড}

\begin{table}[h]\centering\small
\begin{tabular}{ll}
বিষয় & বিবরণ \\
\hline
\textbf{বর্ণনাকারী} & মুহাম্মদ ইবনে কাব আল-কুরাযী (রহ.) থেকে \\
\textbf{গ্রন্থ} & মুসান্নাফ ইবনে আবি শায়বাহ, হাদিস ৩৫৫৯৮; সহীহ বুখারি (তাফসির, সূরা আন-নাবা) ও মুসনাদ আহমাদ, হাদিস ৬২৯ \\
\textbf{অধ্যায়} & মুসান্নাফ — কিতাবুস-সিয়ার/যুহদ; বুখারি — তাফসির \\
\textbf{সনদের মান} & \textbf{সহীহ সনদের আসার} \\
\textbf{মূল বিষয়} & খলিফা এমন শাসন চাইতেন যেখানে মানুষ তাকে নিয়ন্ত্রণ করবে — ক্ষমতার অহংকার নয়, দায়বদ্ধতা \\
\hline
\end{tabular}
\end{table}

\medskip\hrule\medskip

\subsection*{২. পটভূমি (কে, কখন, কোথায়, কেন)}

\textbf{উমার ইবনুল খাত্তাব (রাঃ)}-এর খিলাফতকালে — \lat{power} (ক্ষমতা)-এর শীর্ষ যুগে। মুসলিমদের বিজয়ের পর অনেক এলাকা মুসলিম শাসনের (\lat{rule}) অধীনে চলে আসে। সেখানে নতুন নতুন মানুষের সাথে উমারের \lat{meeting} (সাক্ষাৎ) হতো।

একদিন উমারের কাছে কিছু মানুষ এসেছিল — সম্ভবত বিজিত এলাকার \lat{delegate} (প্রতিনিধি) বা যুবকের দল। উমারের সাথে তাদের কিছু \lat{conversation} (কথা) হচ্ছিল। কথার একপর্যায়ে (\lat{stage}) তাদের একজন উমারকে বলল:

\textbf{\lat{“}আপনার প্রতিনিধিরা আপনার কাছে যা যা (অত্যাচার/অবিচার) করে, তা কি আপনি জানেন না?\lat{”}} — অর্থাৎ শাসনের নিচু স্তরের লোকদের \lat{injustice} (অন্যায়) সম্পর্কে উমারকে \lat{warning} (সতর্ক) করা হলো।

এর জবাবে উমার অত্যন্ত \lat{humble} (বিনয়ী) ভঙ্গিতে এমন একটি কথা বললেন যা ইসলামি শাসনব্যবস্থার \lat{ideal} (আদর্শ) হিসেবে স্মরণীয়।

\medskip\hrule\medskip

\subsection*{৩. পূর্ণ ঘটনার বর্ণনা}

মুহাম্মদ ইবনে কাব আল-কুরাযী (রহ.) বলেন:

\begin{quote}
\textbf{\lat{“}একদল লোক উমার ইবনুল খাত্তাবের (রাঃ) কাছে এল। তাদের একজন বলল: "আপনার প্রতিনিধিরা আপনার কাছে যা কিছু করে, তা কি আপনি জানেন না?"} \\
\textbf{উমার (রাঃ) বললেন: "আল্লাহর কসম! আমি তো চাইছি — যদি আমার মাথার (উপর) ঘাসের একটি মুঠা (হাশিমাহ) থাকে, যার দ্বারা আমার কাছে আল্লাহর সন্তুষ্টি পৌঁছায়, তবে আমি তা নিয়ে (বিনয়ী ও সংযমী) হব।"\lat{”}}
\end{quote}

অর্থাৎ — উমার বললেন: আমি চাই আমার ক্ষমতা এমন একটি ঘাসের মুঠার মতো হোক, যা দিয়ে আল্লাহর সন্তুষ্টি অর্জন করা যায়; ক্ষমতাকে আমি অহংকারের বস্তু নয়, বরং আল্লাহর সন্তুষ্টির মাধ্যম হিসেবে দেখি।

\medskip\hrule\medskip

\subsection*{৪. ধাপে ধাপে বিশ্লেষণ}

\textbf{ধাপ ১ — অভিযোগ/সতর্কতা:} কেউ উমারকে বলল — আপনার \lat{delegate}-রা \lat{injustice} করছে, আপনি কি জানেন না?

\textbf{ধাপ ২ — উমারের সংযমী জবাব:} উমার \lat{anger} (রাগ) করেননি, আত্মপক্ষ সমর্থনের \lat{arrogance} (অহংকার)ও করেননি। বরং তিনি বললেন — আমি তো কামনা করি আমার \lat{power} ঘাসের মুঠার মতোই হোক, যার সাহায্যে আল্লাহর সন্তুষ্টি পাই।

\textbf{ধাপ ৩ — ক্ষমতার প্রকৃত ধারণা:} উমার ক্ষমতাকে নিজের মর্যাদা বা \lat{arrogance}-এর বস্তু মনে করেননি; মনে করতেন \lat{trust} (আমানত) ও \lat{responsibility} (দায়িত্ব) — যার মাধ্যমে আল্লাহর সন্তুষ্টি কামনা করা হয়।

\textbf{ধাপ ৪ — নিজের \lat{weakness} (দুর্বলতা) স্বীকার:} তিনি নিজেকে অপূর্ণ ও ত্রুটিযুক্ত মনে করতেন — প্রতিনিধিদের অন্যায়ের দায়ও তিনি এড়িয়ে যাননি; বরং সংযম ও \lat{humility} (বিনয়)-র ভাষায় জবাব দিয়েছেন।

\textbf{ধাপ ৫ — শাসনের (\lat{leadership}) নীতি:} ক্ষমতায় থাকা মানুষকে \lat{control} (নিয়ন্ত্রণ) করতে হবে — উমার এটাই চেয়েছেন; ক্ষমতার ঘোরে উড়ে যাওয়া নয়।

\medskip\hrule\medskip

\subsection*{৫. শব্দের ব্যাখ্যা}

\begin{table}[h]\centering\small
\begin{tabular}{ll}
শব্দ & অর্থ \\
\hline
\textbf{হাশিমাহ (\arab{حشيمة})} & ঘাসের মুঠা — তৃণের গুচ্ছ \\
\textbf{উমালা (\arab{عمال})} & প্রতিনিধি/কর্মকর্তা — শাসনের নিচু স্তরের নিযুক্ত ব্যক্তিরা \\
\textbf{রিযা (\arab{رضا})} & সন্তুষ্টি — আল্লাহর সন্তুষ্টি \\
\textbf{আমানাহ (\arab{أمانة})} & দায়িত্ব/গুরুপূর্ণ দায় \\
\textbf{তাওয়াদু' (\arab{تواضع})} & বিনয় \\
\hline
\end{tabular}
\end{table}

\textbf{উমারের রূপক:} ঘাসের মুঠা — অতি সামান্য, কোনো মূল্য নেই মানুষের কাছে; কিন্তু আল্লাহর সন্তুষ্টির জন্য তা যদি প্রয়োজন হয়, উমার তাকে ধরে রাখতে চেয়েছেন। ক্ষমতাও তেমনই — তা মানুষের কাছে যত বড়ই হোক, আসল মূল্য আল্লাহর সন্তুষ্টিতে।

\medskip\hrule\medskip

\subsection*{৬. সংশ্লিষ্ট হাদিস ও আয়াত}

\textbf{হাদিস (মুসলিম ৯১):} নবীজি বলেছেন: \textbf{\lat{“}যার হৃদয়ে সরিষার দানার পরিমাণ কিবর থাকবে, সে জান্নাতে প্রবেশ করবে না।\lat{”}}

\textbf{হাদিস (বুখারি ৮৯৩):} নবীজি বলেছেন: \textbf{\lat{“}তোমাদের প্রত্যেকেই শাসক (দায়িত্বশীল) এবং প্রত্যেকেই নিজের প্রজা (অধীনস্থদের) সম্পর্কে জিজ্ঞাসিত হবে।\lat{”}} — ক্ষমতা জিজ্ঞাসার মুখোমুখি।

\textbf{হাদিস (মুসলিম ১৮২২):} নবীজি বলেছেন: \textbf{\lat{“}আল্লাহর কাছে এমন কোনো শাসক নাই যিনি মারা যাবেন — তার আমল সম্পর্কে জিজ্ঞাসিত হবেন না; জান্নাতে যাবেন বা জাহান্নামে।\lat{”}} — শাসকের আমল তীব্র পরীক্ষা।

\textbf{আয়াত (সূরা আন-নাবা ৭৮:১-৩):}
\begin{quote}
\textbf{\lat{“}কী সম্বন্ধে তারা একে অপরকে জিজ্ঞেস করে? মহাসংবাদ সম্পর্কে — যাতে তারা মতভেদ করে।\lat{”}} — আখিরাতের হিসাব।
\end{quote}

\textbf{আয়াত (সূরা আয-যিলযাল ৯৯:৭-৮):}
\begin{quote}
\textbf{\lat{“}অতঃপর যে অণু পরিমাণ সৎকর্ম করবে, তা সে দেখবে; আর যে অণু পরিমাণ অসৎকর্ম করবে, তা-ও সে দেখবে।\lat{”}} — ক্ষমতার হিসাবও এভাবে হবে।
\end{quote}

\medskip\hrule\medskip

\subsection*{৭. গভীর শিক্ষা}

\textbf{১. \lat{power} — \lat{trust}, \lat{arrogance} নয়।}
উমার \lat{power}-কে এমনভাবে দেখেছেন যেন তা ঘাসের মুঠার মতোই ক্ষুদ্র — কারণ প্রকৃত \lat{greatness} (বড়ত্ব) আল্লাহর কাছে। \lat{power}-এ গর্ব করা কিবর।

\textbf{২. শাসকের (\lat{ruler}) নিয়ন্ত্রণ — উম্মাহর অধিকার।}
উমার চেয়েছেন মানুষ তাকে \lat{control} করুক; শাসক নিজেকে জনগণের ঊর্ধ্বে মনে করবেন না — এটাই ইসলামি \lat{principle} (নীতি)।

\textbf{৩. প্রতিনিধিদের জবাবদিহি (\lat{accountability})।}
\lat{complaint} (অভিযোগ) শুনে উমার তদন্তে আগ্রহী — প্রতিনিধিদের অন্যায়ও শাসকের \lat{burden} (দায়) আসে; ভোগ না করে \lat{humble} হয়ে জবাব দিয়েছেন।

\textbf{৪. নিজেকে \lat{arrogance}-মুক্ত রাখা।}
উমার নিজেকে \lat{perfect} (নিখুঁত) ভাবেননি; মানুষকে জিজ্ঞেস করা — প্রতিনিধিদের অন্যায় শোনার প্রস্তুতি — \lat{humility}-র চিহ্ন।

\textbf{৫. আখিরাতের হিসাব (\lat{judgement}) স্মরণ।}
\lat{power}-র প্রতিটি কাজের \lat{judgement} হবে — এই স্মরণই উমারকে \lat{humble} রেখেছে। \lat{power} যেন আখিরাতের হিসাবের কথা ভুলিয়ে না দেয়।

\medskip\hrule\medskip

\subsection*{৮. জীবনে প্রয়োগের পদ্ধতি}

\begin{enumerate}
\item \textbf{\lat{power} বা পদ পেলে:} মনে রাখুন — এটি \lat{trust}; এর জবাব দিতে হবে। তাই \lat{pride} নয়, \lat{accountability} নিয়ে কাজ করুন।
\item \textbf{অধীনদের \lat{criticism} (সমালোচনা) শুনুন:} বাড়ির, প্রতিষ্ঠানের বা সমাজের নেতা হলে \lat{criticism} শুনে \lat{anger} করবেন না — \lat{humble} হয়ে সংশোধন করুন।
\item \textbf{নিজের \lat{weakness} স্বীকার করুন:} নিজেকে \lat{perfect} ভাববেন না; মানুষকে প্রশ্ন করার সুযোগ দিন।
\item \textbf{আল্লাহর সন্তুষ্টিকে লক্ষ্য বানান:} কাজের মধ্যে আল্লাহর \lat{pleasure} (রিযা) খুঁজুন — লোক-দেখানো নয়।
\item \textbf{আখিরাতের হিসাব স্মরণ করুন:} প্রতিটি \lat{decision}-এর আগে ভাবুন — কিয়ামতের দিন এই কাজের জবাব কীভাবে দেব।
\end{enumerate}
\medskip\hrule\medskip

\subsection*{৯. রেফারেন্স}

\begin{itemize}
\item মুসান্নাফ ইবনে আবি শায়বাহ, হাদিস ৩৫৫৯৮ (সহীহ সনদ)
\item সহীহ বুখারি (তাফসির, সূরা আন-নাবা)
\item মুসনাদ আহমাদ, হাদিস ৬২৯
\item সহীহ মুসলিম, হাদিস ৯১ ও ১৮২২; সহীহ বুখারি, হাদিস ৮৯৩
\item সূরা আন-নাবা ৭৮:১-৩; সূরা আয-যিলযাল ৯৯:৭-৮
\end{itemize}

\end{document}
